% !TeX program = latexmk
% Thesis or personal report (BEP, master thesis, internship report): EE report
% cover, title page, abstract, contents, chapters in the order of a scientific
% argument (SM1-SM6, TU/e thesis template), references, appendices.
% Build: make (or latexmk -pdf thesis.tex). Style: tuethesis.sty.
% References: references.bib. Writing rules: WRITING.md in the LaTeX templates.
% Printed and bound: add twoside to the class options (chapters start on the
% right, back cover on an even page).
\documentclass[a4paper, 11pt]{report}
\usepackage{tuethesis}   % [font=noto] Noto Sans, [font=cm] Computer Modern
\addbibresource{references.bib}

\reporttype{Master thesis}               % Bachelor thesis, Internship report, ...
\reporttitle{Title of the thesis}
\reportsubtitle{}                         % empty: no subtitle
\reportauthors{Name Surname}
\reportkeywords{}                          % PDF keywords, comma-separated
\reportstudentid{1234567}                 % empty: hidden
\reportversion{v0.1}                      % v0.x drafts, v1.0 first release
\reportgroup{}                            % research group, empty: hidden
\reportdepartment{Department of Electrical Engineering}
\reportsupervisor{%
  First committee member (supervisor)\\
  Second committee member\\
}
% \reportcoverimage{figures/cover.jpg}    % dark photo; none: plain TU/e red
% \coverdarkentrue                        % darken the photo bottom if needed
% \reportdate{25 November 2026}           % default: \today
% \accentcolor{HTML}{003399}              % secondary color for headings etc.; default Scarlet

\begin{document}

\makecover
\maketitlepage

\frontmatter
% Abstract: one sentence (two at most) for each point, written last.
% 1. The topic, in words your reader understands.
% 2. The problem: the key research question.
% 3. Why nobody has answered it yet (the gap in the literature).
% 4. How you tackled it: the big new idea.
% 5. How you did the research (experiments, software, case studies).
% 6. The key impact: what it means and why others should care.
% No citations.
\frontchapter{Abstract}
State the topic, the research question, the gap, the idea, the method, and the
impact, one sentence each.

\tableofcontents
\listoffigures
\listoftables

\frontchapter{List of abbreviations}
\begin{acronym}[ADC]
  \acro{ADC}{analog-to-digital converter}
\end{acronym}

\mainmatter
% ------------------------------------------------------------------------------
\chapter{Introduction}
\label{ch:introduction}
% The whole argument in one chapter, readable on its own by anyone in the
% field: problem, method, relevant execution details, results, and their
% interpretation. Write it last; keep notes for each section while writing the
% later chapters.
% Opening paragraph: the chapter's topic, its sections in order with forward
% references.
This chapter summarizes the work.
\Cref{sec:context} introduces the topic, \cref{sec:state-of-the-art} the state
of the art, \cref{sec:research-question} the research question,
\cref{sec:method} the method, and \cref{sec:findings} the findings.

\section{Context and topic (SM1)}
\label{sec:context}
% The general subject, the specific topic of interest, and what society or
% science wants to know about it. Understandable for a general scientific
% audience.

\section{State of the art (SM1)}
\label{sec:state-of-the-art}
% Established facts, methods and results on the topic, ending with the
% knowledge gap this work addresses.

\section{Research question (SM2)}
\label{sec:research-question}
% The specific part of the gap this work addresses, with its assumptions made
% explicit.

\section{Method or approach (SM3, SM4)}
\label{sec:method}
% The steps someone else needs to reproduce the work, including how success is
% evaluated. Refer to the chapters with the details.

\section{Findings (SM5, SM6)}
\label{sec:findings}
% The evaluation setup, the results, and their interpretation: a summary of
% \cref{ch:evaluation}.

% ------------------------------------------------------------------------------
\chapter{Background (SM1)}
\label{ch:background}
% Exactly the preliminaries a student in your program who did not study this
% subject needs to follow the rest; related work that the state of the art in
% the introduction only summarized.

% ------------------------------------------------------------------------------
\chapter{Problem exposition (SM1--SM3)}
\label{ch:problem}
% Optional: only when the introduction does not give enough detail about the
% problem context, the data, the requirements, or the method.

\section{Context (SM1)}
\section{Detailed research questions (SM2)}
\section{Detailed method (SM3)}

% ------------------------------------------------------------------------------
\chapter{First sub-problem}
\label{ch:first}
% From here on: execution (SM4), results (SM5) and interpretation (SM6), one
% sub-problem per chapter or section, each again as a small scientific
% argument. First paragraph: the sub-problem and the requirements a solution
% must meet, based on the preceding chapters.
The spring force follows Hooke's law, \cref{eq:hooke}.
\begin{equation}
  F = k \, \delta
  \label{eq:hooke}
\end{equation}
With $k = \SI{250}{\newton\per\metre}$ and a compression of \SI{40}{\milli\metre}, the force is \SI{10}{\newton}~\cite{knuth1986}.
The \ac{ADC} samples the force sensor.

\begin{figure}[htbp]
  \centering
  \includegraphics[width=0.5\textwidth]{example-image}
  \caption{Every figure is referred to and explained in the text.}
  \label{fig:example}
\end{figure}

\section{Design}
\label{sec:design}
\Cref{fig:example} shows the design.

% ------------------------------------------------------------------------------
\chapter{Second sub-problem}
\label{ch:second}

% ------------------------------------------------------------------------------
\chapter{Evaluation}
\label{ch:evaluation}
% The most formal chapter: the evaluation of the whole work against the
% research question.

\section{Objective (SM2)}
\section{Setup (SM3)}
\section{Execution (SM4)}
\section{Results (SM5)}
% Describe and summarize the measurements; do not interpret them here.
\section{Discussion (SM6)}
% Interpret the results: what they mean for the research question, and their
% limitations.

% ------------------------------------------------------------------------------
\chapter{Conclusion}
\label{ch:conclusion}
% The problem, the solution, the results that answer the research question, and
% future work. No citations.

% ------------------------------------------------------------------------------
\printbibliography[heading=bibintoc, title={References}]

\appendix
\chapter{First appendix}
\label{app:first}
Long derivations and raw data, referred to from the text as \cref{app:first}.

\makebackcover
\end{document}
